\chapter{The automated hybrid surrogate-MCMC pipeline}
\label{ch:hybrid}

The procedure of \cref{ch:surrogate} works, but leaves every decision to the user and gives no
guarantee at the end. \code{slurm_mcmc_hybrid} automates it: it iterates fit, sample and validate
until measured criteria say the surrogate posterior can be trusted, recycling every expensive
evaluation it makes along the way.

\section{The loop}
\label{sec:hybrid-loop}

\begin{enumerate}[leftmargin=*, label=\textbf{\arabic*.}]
  \item \textbf{Round 0.} A \emph{short} expensive MCMC, plus optional
        \code{num_regularization_points}, generates the initial training set. Its walkers --- the
        parallel MCMC chains of \cref{ch:mcmc} --- are evaluated in parallel on the cluster, so
        their number is this stage's parallelism. The regularization points are the only
        evaluations in the whole run placed without reference to a posterior, and so the only ones
        that can find mass the surrogate does not already know about.
  \item \textbf{Fit} a fast surrogate: a Gaussian process by default (\cref{sec:hybrid-surrogates}).
  \item \textbf{Sample} the surrogate with a cheap, vectorized MCMC that needs no cluster, confined
        to \code{param_bounds} or, by default, to an envelope grown from the training data
        (\cref{sec:unbounded}). A fresh sampler is built every round, because each round samples a
        different surrogate, with its own deliberately small set of walkers. This chain is kept
        cheap, because its posterior is discarded at the next fit.
  \item \textbf{Validate.} Evaluate the expensive function, in parallel, on
        \code{num_validation_points} drawn from the surrogate posterior, and score the criteria of
        \cref{sec:hybrid-measured}.
  \item \textbf{Refine.} If the criteria are not met, the validation evaluations --- which sit
        exactly where the surrogate needs improving, and are already paid for --- join the training
        set, and the loop returns to step 2. Optionally, each round also runs a few expensive MCMC
        steps started from the current surrogate posterior (\code{num_expensive_iters_per_round}).
  \item \textbf{Verify.} A pass reached on a cheap chain is only provisional. The same surrogate is
        re-sampled to reporting depth (\cref{sec:two-depths}) and one more validation batch is
        spent on \emph{that} posterior. If it still passes, it is the posterior returned; if not,
        refinement resumes. No evaluation is wasted either way: the refinement batch joins the
        training set as soon as verification starts, and a rejected verification's batch joins it
        too.
\end{enumerate}

Verification buys two things. The returned posterior is sampled properly, rather than being
whatever a deliberately cheap chain happened to produce. And the decision to stop rests on two
independent measurements rather than one --- which matters, because a single measurement from 100
validation batch is noisy: two batches scored against the \emph{same} surrogate can land on opposite
sides of the threshold.

\begin{keypoint}[Starting from existing evaluations]
If expensive evaluations already exist --- from an ordinary \code{slurm_mcmc} run, a parameter scan,
an earlier study --- pass them as \code{initial_train_points} and \code{initial_train_values}. Round 0
is skipped and the first surrogate is fitted on them directly; \code{init_points} is then not needed,
and the evaluations do not count towards \code{num_expensive_evals}, which records only what the run
itself spent. The regularization points are part of round 0, so they are skipped too, with a warning.
\begin{lstlisting}
pool = status['slurm_pool']                 # from a previous slurm_mcmc
result = slurm_mcmc_hybrid(..., initial_train_points=pool.points_history,
                           initial_train_values=pool.values_history[:, 0])
\end{lstlisting}
\end{keypoint}

\section{What is measured}
\label{sec:hybrid-measured}

On the validation points $\theta_i$, write the log-ratio of the expensive to the surrogate
log-probability, and the self-normalised importance weights of \cref{sec:importance}:
\begin{equation}
  d_i = \log\pexp(\theta_i) - \log\psur(\theta_i),
  \qquad w_i = e^{d_i},
  \qquad \tilde w_i = \frac{w_i}{\sum_j w_j}.
\end{equation}
The weights are the \emph{ratio}, not the expensive density, because the points were drawn from the
surrogate posterior and so already arrive with density $\propto \psur$: the sampling density has to
be divided out before the target density is multiplied in.

\begin{keypoint}[Nats]
Every logarithm here is natural, and the unit of a log-probability is the \emph{nat} --- which is
what ``[nats]'' means on the figures below. A difference of $\Delta$ nats in log-probability is a
factor $e^{\Delta}$ in probability: 1~nat is a factor of 2.7, 0.1~nats about 10\%, and 0.01~nats
about 1\%.
\end{keypoint}

\subsection{Surrogate accuracy}

\begin{equation}
  \bar d = \sum_i \tilde w_i\, d_i,
  \qquad
  \varepsilon = \sqrt{\sum_i \tilde w_i\, (d_i - \bar d)^2}\quad[\text{nats}].
  \label{eq:epsilon}
\end{equation}
$\varepsilon = 0.05$ means probability ratios are accurate to about 5\%. The mean $\bar d$ is
subtracted because a log-probability is only defined up to a constant, so only the \emph{spread} of
$d$ is a real error. The weighting makes $\varepsilon$ an average under the \emph{expensive}
posterior: an error where the true probability is high counts far more than the same error in the
tails.

\subsection{Degeneracy guard}

$\varepsilon$ alone is not safe. If the weights collapse onto a few points, it is computed from
those points alone and can be arbitrarily small --- exactly zero for a single dominating point ---
while the surrogate is in fact terrible. The number of validation points that effectively carry the
weight is
\begin{equation}
  \ESSw = \frac{\left(\sum_i w_i\right)^2}{\sum_i w_i^2} \in [1, N],
\end{equation}
and a run is never declared converged while $\ESSw <$ \code{min_ess_weights}. $\ESSw$ counts
validation points and is unrelated to the ESS of a chain, $N/\tau$, of \cref{eq:ess}.

\subsection{Posterior stability}

Surrogate accuracy and posterior convergence are different questions: a perfect surrogate still
leaves a drifting posterior if its chain is too short. Between rounds, the change in each
parameter's posterior mean $m_j$ and standard deviation $s_j$ is measured in units of that
parameter's posterior width, and the largest is kept:
\begin{equation}
  \delta = \max_j
  \frac{\max\!\left(\,\big|m_j^{\mathrm{new}} - m_j^{\mathrm{old}}\big|,\ \big|s_j^{\mathrm{new}} - s_j^{\mathrm{old}}\big|\,\right)}
       {\tfrac12\left(s_j^{\mathrm{new}} + s_j^{\mathrm{old}}\right)} .
  \label{eq:delta}
\end{equation}
$\delta = 0.1$ means nothing moved by more than a tenth of a posterior width --- a statement about
the answer rather than about the chain. Two estimates of the \emph{same} posterior still differ by
chance, by roughly $\sqrt{1/\ESS_{\mathrm{new}} + 1/\ESS_{\mathrm{old}}}$ widths, where ESS is the
chain's and ``old'' is the previous round's reported posterior. The tolerance actually applied
therefore never drops below three times that noise floor:
\begin{equation}
  \delta \;\le\; \max\!\left(\texttt{posterior\_shift\_tolerance},\;
  3\sqrt{\frac{1}{\ESS_{\mathrm{new}}} + \frac{1}{\ESS_{\mathrm{old}}}}\right).
  \label{eq:stability}
\end{equation}
Without the floor a tolerance of 0.1 would be unreachable at realistic chain depths: at the ESS of
95--380 in the worked example, the floor alone is 0.23--0.42. The test needs two rounds, so a run
cannot converge in round 1 while it is enabled; \code{posterior_shift_tolerance=None} switches it
off.

\subsection{The convergence test}

All three conditions must hold simultaneously:
\begin{gather}
  \varepsilon \le \texttt{log\_error\_threshold},
  \qquad
  \ESSw \ge \texttt{min\_ess\_weights},
  \notag\\
  \delta \le \max\!\left(\texttt{posterior\_shift\_tolerance},\ \text{noise floor}\right).
\end{gather}

\begin{cautionbox}[What the test cannot see]
Every statistic is evaluated on points drawn from the \emph{surrogate} posterior. The test therefore
detects a surrogate posterior that is too \textbf{broad}, and cannot detect one that is too
\textbf{narrow}: mass the surrogate never finds is never sampled, never scored, and never enters any
statistic. The worked example below misses a tail in exactly this way. If the target may be
hierarchical or multimodal, spend some budget on \code{num_regularization_points}, disperse
\code{init_points}, and consider an independent short expensive chain as a check.
\end{cautionbox}

\section{Two sampling depths}
\label{sec:two-depths}

The two chain criteria of \cref{sec:mcmc-convergence} are applied at two depths. The surrogate
chain runs \code{num_surrogate_iters} iterations, and is extended in blocks of that size, on the same
sampler, until the criteria of its tier hold --- or \code{max_surrogate_iters_multiplier} blocks
have been spent, which is logged as a warning.

\begin{center}
\small
\begin{tabularx}{\linewidth}{@{}l L L@{}}
  \toprule
  & \textbf{refining} --- posterior discarded at the next fit & \textbf{verifying} --- posterior reported \\
  \midrule
  effective samples, $N/\tau$ & \code{refine_surrogate_ess}: \code{None} $\to 10\,d$ & \code{final_surrogate_ess}: \code{None} $\to 100\,d$ \\
  chain length, $\Lc/\tau$    & \code{refine_iters_per_tau}: \code{None} $\to$ not demanded & \code{final_iters_per_tau}: $50$ \\
  \bottomrule
\end{tabularx}
\end{center}

The defaults are the two ends of the advised band of $10$--$100\times d$. The length criterion is not
demanded while refining, for two reasons: that posterior is discarded, and $\tau$ cannot be measured
reliably on a short chain anyway --- that is what $\Lc \ge 50\tau$ \emph{means} --- so a floor on it
would be enforced against a number that is itself unreliable. Samples and $\tau$ are always taken
after discarding \code{surrogate_burnin_fraction} of the chain.

\code{max_verification_attempts}, by default \code{None}, lets a failed verification send the run
back to refining as often as the target needs, since \code{num_rounds_max} bounds it anyway. A finite
cap means \emph{give up after this many rejections} and ends the run: a run that can no longer
verify can never converge. \code{0} switches verification off.

\section{Surrogate models}
\label{sec:hybrid-surrogates}

\code{surrogate} accepts one of the following, or any object with \code{fit(X, y)} and
\code{predict(X)} --- and \code{predict_std(X)} too, if the uncertainty penalty should be active:

\begin{itemize}
  \item \textbf{\code{'gp'}} (default), a \code{GaussianProcessSurrogate} --- the recommended choice:
        accurate, smooth, and it reports its own uncertainty.
  \item \textbf{\code{'polynomial'}}, a \code{PolynomialSurrogate} --- a cheap fallback for
        posteriors that really are low-order. It extrapolates badly otherwise, and has no
        uncertainty, so the uncertainty penalty is inactive with it.
  \item \textbf{\code{DistributedGPSurrogate(...)}} --- an ensemble of Gaussian processes, for when
        the training set has grown too large for a single one to be refitted every round.
\end{itemize}

The first two wrap standard scikit-learn models:
\code{GaussianProcessRegressor} for the Gaussian process, and a ridge regression on polynomial
features for the polynomial. The Gaussian process first fits a polynomial mean function, so that the
process only has to model what the polynomial misses. The third,
\code{DistributedGPSurrogate}, is not a standard model, and is described in \cref{sec:rbcm}.

\subsection{The uncertainty penalty}

During the surrogate MCMC the surrogate is evaluated pessimistically, as
\begin{equation}
  \log\psur(\theta) = \mu(\theta) - \kappa\, \sigma(\theta),
\end{equation}
where $\mu$ and $\sigma$ are its prediction and uncertainty and $\kappa$ is
\code{surrogate_uncertainty_penalty}. A Gaussian process's uncertainty grows away from its training
data, so walkers are turned back from regions where the surrogate would be extrapolating.

\subsection{DistributedGPSurrogate}
\label{sec:rbcm}

\textbf{Why it exists.} The surrogate is refitted every round on the whole training set, and the
cost of fitting a Gaussian process grows steeply --- roughly with the cube of the number of points.
Early in a run that is negligible, but every round adds evaluations, so on a long run the fit comes
to dominate and eventually stops being affordable at all. \code{DistributedGPSurrogate} is the way
past that point.

\textbf{What it does.} It is a robust Bayesian Committee Machine (Deisenroth \& Ng 2015, with Cao \&
Fleet's 2014 weights). It splits the training set among independent Gaussian-process
\emph{experts} and recombines their predictions in closed form. It attacks cost by parallelism
rather than approximation: the expert fits are independent, so on a cluster the wall clock is that of
a single expert.

At any point, every expert predicts a value and an uncertainty, and the ensemble takes a weighted
average of the values. An expert's weight is how much it has learned at that point: how far its
uncertainty there has fallen below the uncertainty it had before seeing any data. An expert sitting
on its own training points is confident and dominates; an expert far from them is back where it
started and gets no weight at all. The weights add up to one, so a single expert reproduces itself
exactly --- a one-expert ensemble \emph{is} the dense Gaussian process --- and where no expert has
learned anything, the ensemble falls back to the experts' starting uncertainty, which is large. That
large uncertainty is what keeps the uncertainty penalty working away from the data. Each expert is
judged against the uncertainty \emph{it} started with: experts routinely settle on scales far from
the nominal one, and judging them against it would make every expert look uninformed away from its
data.

\begin{argtable}
  \code{max_points_per_expert} & \code{None} & set this for large training sets: it grows the \emph{number} of experts with the data instead of the size of each, which keeps fitting affordable \\
  \code{num_experts} & \code{8} & the number of experts, when \code{max_points_per_expert} is not set \\
  \code{min_points_per_expert} & \code{200} & below this many points per expert, fewer experts are used, down to one; an expert fitted on a few dozen points cannot identify its own kernel \\
  \code{partition} & \code{'kmeans'} & how points are split: \code{'kmeans'} (spatial clusters) or \code{'random'} \\
  \code{n_restarts_optimizer} & \code{2} & restarts of each expert's hyperparameter fit; more give better fits for little cost \\
  \code{parallel} & \code{'none'} & \code{'joblib'}: one process per core; \code{'slurm'}: one job array across nodes, worth it only when a single expert fit is long compared with the queue wait \\
  \code{n_jobs}, \code{cluster}, \code{submitit_kwargs} & \code{-1}, \code{'slurm'}, \code{None} & settings of the \code{'joblib'} and \code{'slurm'} modes \\
  \code{work_dir} & \code{None} & scratch location for \code{'slurm'}, \code{work_dir/surrogates/dgp_experts}, removed once the experts are back; inherited from the pipeline's, and must be on a shared filesystem \\
  \code{job_name} & \code{None} & for \code{'slurm'}, the expert job arrays are named \code{job_name_dgp_fit<n>}; inherited from the pipeline's \\
  \code{max_fit_retries} & \code{2} & failed expert jobs are retried individually, then fitted locally \\
  \code{trend_degree} & \code{2} & the polynomial mean function, as for the dense Gaussian process \\
\end{argtable}

Two behaviours are worth knowing. The prediction cannot be computed without the uncertainties,
because the weights come from them, so turning the uncertainty penalty off saves the ensemble
nothing. And cost moves rather than disappears: every MCMC step queries all the experts, so fitting
becomes nearly free while sampling becomes the dominant cost, and end to end the ensemble can be
slower than the dense Gaussian process even though it fits far faster. It pays once the dense fit
itself is the bottleneck.

\section{Training data}

The training set only grows: every expensive evaluation is kept. What enters each fit can be
limited. \code{train_log_prob_trim_range} drops points more than that many nats below the best
one --- their probability is $e^{-100}$ at the default, so no MCMC will ever visit them, and they only
make the fit more expensive --- as long as at least \code{min_train_points_after_trim} points remain.
Below about 20 nats accuracy collapses, because the surrogate then has no evidence that the
probability is low away from the peak. \code{max_train_points} caps the fit with a uniform random
subsample, which preserves spatial coverage; keeping the highest-probability points instead would
discard exactly the far-field points that stop the surrogate extrapolating. \code{train_log_prob_floor}
drops points below an absolute value, such as penalty or failure values.

\section{Running without bounds}
\label{sec:unbounded}

\code{param_bounds} defaults to \code{None}, and then no box is imposed. A box cannot simply be
dropped, though, because a surrogate says nothing reliable away from its training data, and both ways
it can be wrong there are fatal to an MCMC: with a fitted quadratic trend the surrogate tends to that
quadratic, which runs off to $+\infty$ along any direction of positive curvature; without a trend it
tends to a constant, which is improper over infinite volume. Either way the chain walks away and
never returns.

So the training set's own geometry decides where the surrogate may be believed. Let $r(\theta)$ be
the Mahalanobis radius of $\theta$ under the training points' mean and covariance, and $R$ the largest
such radius among them. A quartic penalty switches on beyond $2R$, and the log-probability is
$-\infty$ beyond $10R$, which makes the posterior proper outright. Nothing acts inside $2R$, so the
posterior over the data is untouched; and because both radii are multiples of $R$, they grow as
evaluations accumulate outwards, so a genuinely long tail is followed round after round.

With the default Gaussian process the envelope never actually fires: its uncertainty grows away from
the data, so the uncertainty penalty turns the walkers back long before the envelope --- on the
worked example, the largest radius the walkers reached was 3.51 with the envelope switched on and off
alike. It is load-bearing only for surrogates without an uncertainty, such as the polynomial.

\needspace{22\baselineskip}
\section{Usage}
\label{sec:hybrid-usage}

\begin{lstlisting}
import numpy as np
from slurmcmc.hybrid import slurm_mcmc_hybrid

def log_prob(theta):                            # the expensive log-probability
    ...

num_params = 3
init_points = np.random.uniform(-2, 2, (10 * num_params, num_params))   # round-0 walkers

result = slurm_mcmc_hybrid(log_prob_fun=log_prob, init_points=init_points,
                           num_expensive_iters=10, num_regularization_points=50,
                           log_error_threshold=0.01, num_rounds_max=30,
                           cluster='slurm', work_dir='hybrid',
                           save_restart=True, random_seed=0)

print(result['converged'], result['num_rounds'], result['num_expensive_evals'])
samples = result['samples']                     # the posterior, as flat draws
\end{lstlisting}

Most arguments can stay at their defaults. The ones that decide a run are its round-0 budget ---
the number of walkers in \code{init_points}, \code{num_expensive_iters} and
\code{num_regularization_points} --- the accuracy it must reach, \code{log_error_threshold}, and
how many rounds it may take, \code{num_rounds_max}. To start from evaluations that already exist,
pass \code{initial_train_points} and \code{initial_train_values} instead of \code{init_points}
(\cref{sec:hybrid-loop}); to use another surrogate, pass \code{surrogate}
(\cref{sec:hybrid-surrogates}). With \code{remote=True} the call returns a submitit job, and
\code{result = job.result()} gives the same dictionary (\cref{sec:remote}).

\code{example_mcmc_hybrid.py} spells out every argument on its own line and is the best starting
point for a real problem; \cref{sec:hybrid-args} describes them all, and the returned dictionary is
described at the end of this chapter.

\section{Worked example: a 3D Rosenbrock}
\label{sec:hybrid-example}

This is exactly what \code{example_mcmc_hybrid.py} does, so every number below is reproducible by
running it. The target is minus the 3D Rosenbrock function of \cref{eq:rosenbrock}: a narrow,
curved, banana-shaped posterior. Because it is analytic, a plain expensive MCMC on it is affordable,
and that gives a reference to judge the hybrid run against --- so the reference comes first.

\subsection{The expensive reference}

The reference is a plain \code{slurm_mcmc} on the same target, with no surrogate anywhere: 100\,000
iterations with 30 walkers, \textbf{3\,000\,000 evaluations}. It is cached separately from the
pipeline's results, so changing a hybrid setting does not pay for it again. A reference is only a
yardstick once it has converged itself, so the example draws its posterior at five points along the
chain, after discarding the first 20\% as burn-in.

\begin{figure}[H]
  \centering
  \includegraphics[width=0.5\linewidth]{example_mcmc_hybrid_reference_convergence.png}
  \caption{The reference posterior at five stopping points along the chain, with $\Lc/\tau$ and
           the effective sample size for each in the legend. The figures label the effective sample
           size ESS/$\tau$, for the number of samples divided by $\tau$; $\tau$ is the largest over
           the parameters.}
  \label{fig:hybrid-reference}
\end{figure}

\begin{center}
\small
\begin{tabular}{@{}rrlrr@{}}
  \toprule
  \textbf{stop (iteration)} & $\Lc$ & $\tau$ \textbf{per parameter} & $\Lc/\tau$ & \textbf{ESS} \\
  \midrule
  36\,000  & 16\,000 & 477, 441, 473    & 34 & 1\,006 \\
  52\,000  & 32\,000 & 645, 481, 529    & 50 & 1\,488 \\
  68\,000  & 48\,000 & 754, 533, 617    & 64 & 1\,910 \\
  84\,000  & 64\,000 & 828, 560, 649    & 77 & 2\,318 \\
  100\,000 & 80\,000 & 1\,035, 589, 685 & 77 & 2\,320 \\
  \bottomrule
\end{tabular}
\end{center}

The distributions agree from the third stop on, and $\Lc/\tau$ clears 50. But the caveat of
\cref{sec:mcmc-convergence} applies: $\Lc/\tau$ only means something once $\tau$ has stopped
growing, and $\tau(x_0)$ has not --- its last step is the largest. The cause is a long arm of the
posterior at negative $x_0$: walkers enter it rarely (about 8\% of the mass lies at $x_0 < 0$), and a
single visit can last over 13\,000 iterations, so the longer the chain, the longer the excursions it
has seen. The moments are stable within their error bars, so the comparison below stands, but the
reference's own uncertainty on the width in $x_0$ may be understated.

\subsection{The hybrid run}

A Gaussian-process surrogate with the default quadratic mean function. The settings that differ from
the library defaults:

\begin{center}
\small
\begin{tabularx}{\linewidth}{@{}>{\raggedright\arraybackslash}p{0.40\linewidth}>{\raggedright\arraybackslash}p{0.12\linewidth}L@{}}
  \toprule
  \textbf{setting} & \textbf{value} & \textbf{why} \\
  \midrule
  expensive walkers, \code{len(init_points)} & $10\,d = 30$ & the cluster parallelism of the expensive stage \\
  \code{num_expensive_iters}, \code{num_regularization_points} & 10, 50 & a deliberately small round-0 budget, so refinement has work to do \\
  \code{num_surrogate_iters}, \code{max_surrogate_iters_multiplier} & 5\,000, 30 & room for the verification chain to reach reporting depth \\
  \code{log_error_threshold} & 0.01 & demanding, so the loop refines for several rounds \\
  \code{posterior_shift_tolerance} & \code{None} & stability not tested: the run stops on accuracy alone \\
  \code{num_rounds_max} & 30 & this target needs more than a handful \\
  \code{save_posterior_samples} & 2\,000 & keeps each round's posterior \\
  \bottomrule
\end{tabularx}
\end{center}

\begin{keypoint}[Result]
Converged in \textbf{7 rounds on 1\,380 expensive evaluations}, about 16 minutes --- 0.05\% of the
reference's evaluations. Three verification pushes were spent; two were rejected.
\end{keypoint}

\subsection{Convergence, round by round}

The posterior shift is in units of \textbf{posterior standard deviations} --- the ``[sd]'' on its
axis and in the table below. It is $\delta$ of \cref{eq:delta}: a shift of 0.1 means that, since the
previous round, no parameter's posterior mean or standard deviation moved by more than a tenth of that
parameter's posterior standard deviation.

A verified round has \textbf{two rows}: $n$ for the refinement chain, whose metrics passed and so
triggered the verification, and $n$v for the verification chain that then judged it. $n_{\mathrm{exp}}$
counts the expensive evaluations accumulated before the round, and $n_{\mathrm{train}}$ those that
entered the fit; the two rows share one fit, so both are given on the first row only.

\begin{figure}[H]
  \centering
  \includegraphics[width=\linewidth]{example_mcmc_hybrid_convergence.png}
  \caption{Top: the criteria, surrogate accuracy $\varepsilon$ and posterior shift $\delta$, with
           the region a converged run must stay out of shaded. Bottom: how deeply the surrogate
           chain was sampled, $\Lc/\tau$ and ESS against their floors. Green squares are
           verification chains, blue circles refinement chains.}
  \label{fig:hybrid-convergence}
\end{figure}

\begin{center}
\small
\begin{tabular}{@{}lrrrrrrr@{}}
  \toprule
  \textbf{round} & $n_{\mathrm{exp}}$ & $n_{\mathrm{train}}$ & $\varepsilon$ [nats] & $\delta$ [sd] & $\Lc$ & $\Lc/\tau$ & ESS \\
  \midrule
  1  & 380    & 123 & 0.7636 & ---   & 5\,000  & 20.3 & 130 \\
  2  & 480    & 223 & 0.0179 & 0.867 & 5\,000  & 18.9 & 121 \\
  3  & 580    & 323 & 0.0078 & 0.043 & 5\,000  & 19.0 & 122 \\
  3v &        &     & 0.0239 & 0.153 & 20\,000 & 51.6 & 330 \\
  4  & 780    & 523 & 0.0098 & 0.224 & 5\,000  & 18.4 & 118 \\
  4v &        &     & 0.0264 & 0.149 & 40\,000 & 59.4 & 380 \\
  5  & 980    & 723 & 0.0280 & 0.101 & 5\,000  & 17.1 & 109 \\
  6  & 1\,080 & 823 & 0.0102 & 0.243 & 5\,000  & 18.1 & 116 \\
  7  & 1\,180 & 923 & 0.0035 & 0.117 & 5\,000  & 14.8 & 95 \\
  7v &        &     & 0.0036 & 0.112 & 30\,000 & 50.6 & 324 \\
  \bottomrule
\end{tabular}
\end{center}

\begin{itemize}
  \item \textbf{Verification did the deciding work.} Rounds 3 and 4 passed on their refinement
        chains, at $\varepsilon = 0.0078$ and $0.0098$; the \emph{same surrogates}, re-sampled
        properly and scored on fresh batches, gave $0.0239$ and $0.0264$ and were rejected. Round 7
        passed both, and rounds 5 and 6 missed the threshold on their refinement chains. That the
        fresh estimate came out higher each time is expected: a round is only verified when its
        refinement $\varepsilon$ happened to land below the threshold, and a single 100-point
        estimate scatters. A stop decided on one batch is partly luck.
  \item \textbf{The two sampling depths are explicit in the bottom row.} Refinement chains ran the
        minimum 5\,000 iterations at $\Lc/\tau \approx 15$--20, deliberately short of 50, with ESS of
        95--130 against their floor of $10\,d = 30$. The three verification chains were extended to
        20\,000--40\,000 iterations, until both $\Lc/\tau \ge 50$ and ESS $\ge 100\,d = 300$ held.
  \item \textbf{Stability was not tested, so the shift column is for information.} From round 3 on,
        every shift is 0.04--0.24 posterior widths, against a noise floor of 0.23--0.42. Every one
        sits below its floor, so had the default tolerance of 0.1 been on, the run would have taken
        exactly the same path.
  \item \textbf{$n_{\mathrm{exp}}$ against $n_{\mathrm{train}}$ is the trim at work.} The yield
        $n_{\mathrm{train}}/n_{\mathrm{exp}}$ climbs from 32\% in round 1 to 78\% by round 7. Every
        trimmed point comes from round 0: at the last fit, 207 of the 330 points of the short initial
        chain, which set out from dispersed walkers, and all 50 regularization points. Not one
        of the 800 points added after round 0 was trimmed, because they were drawn from the posterior
        itself.
  \item \textbf{A verified round costs 200 evaluations, not 100}, because both the refinement batch
        and the verification batch are recycled into the training set.
\end{itemize}

\subsection{The posterior, round by round}

\begin{figure}[H]
  \centering
  \includegraphics[width=0.82\linewidth]{example_mcmc_hybrid_hybrid_posterior_per_round.png}
  \caption{Each round's posterior as the pipeline left it --- for a verified round, the verification
           chain's --- one colour per round, with the expensive reference dashed in black.}
  \label{fig:hybrid-rounds}
\end{figure}

Round 1 ($\varepsilon = 0.76$) is visibly wrong: a single narrow peak, too high in all three
parameters. From round 2 on the rounds follow the reference's shape, and they reach further into the
long arm at negative $x_0$ as the training data grows: the mass at $x_0 < -0.5$ is 0.6\% in round 2,
1.4\% in round 3, and between 1\% and 4\% after that, against the reference's 2.8\%. What no round
reaches is the other end of the banana, the upper tails of $x_1$ and $x_2$ where the valley climbs
steeply. $\varepsilon$, being measured on points drawn from the surrogate posterior, cannot see a
region the chain never visits. The samples behind the figure are kept in
\code{work_dir/posteriors/round{k}.npy} when \code{save_posterior_samples} is set; nothing else
records them.

\subsection{Where the evaluations went}

\begin{figure}[H]
  \centering
  \includegraphics[width=0.96\linewidth]{example_mcmc_hybrid_log_prob_maps.png}
  \caption{The expensive log-probability in each parameter pair, with the remaining parameter held at
           the optimum, and 200 draws each from the hybrid posterior (red) and the reference (light
           blue). The colour scale is in nats, clipped 200 below the peak.}
\end{figure}

The two layers live in different spaces: the background is a \emph{slice} with the remaining
parameter pinned at the optimum, while the dots are a \emph{projection} of the full 3D posterior. A
bright patch with no dots on it, such as the lobe near $x_0 = -1$, is therefore not by itself evidence
of a missed mode. What does matter is visible in the top-left panel: the red draws follow the valley
along the arm at negative $x_0$ down to about $-0.8$, so this run did find it, but at the other end
they stop near $x_1 \approx 2.1$, where the reference continues higher.

\subsection{The posterior against the reference}

\begin{figure}[H]
  \centering
  \includegraphics[width=0.82\linewidth]{example_mcmc_hybrid_posterior.png}
  \caption{The hybrid posterior (black) against the expensive reference (blue).}
\end{figure}

\begin{center}
\small
\begin{tabular}{@{}lccc@{}}
  \toprule
  & $x_0$ & $x_1$ & $x_2$ \\
  \midrule
  hybrid standard deviation    & 0.492 & 0.527 & 0.962 \\
  reference standard deviation & $0.493 \pm 0.016$ & $0.594 \pm 0.006$ & $1.296 \pm 0.030$ \\
  hybrid too narrow by         & 0\% & \textbf{11\%} & \textbf{26\%} \\
  \bottomrule
\end{tabular}
\end{center}

The reference's error bars are batch means over ten blocks of 8\,000 iterations. The centres agree
to within 0.14 posterior widths. The width in $x_0$ agrees too: the arm with $x_0 < -0.5$ holds
3.8\% of the hybrid's mass against $2.8\% \pm 0.5\%$ of the reference's. The widths in $x_1$ and
$x_2$ do not, and the upper tails say why: $x_1 > 2$ holds 4.0\% of the reference's mass and 0.4\%
of the hybrid's, and $x_2 > 4$ the same. The run never sampled that end of the banana, so it never
scored it, and every criterion passed regardless. Which tail is cut is partly chance: an earlier run
of this example, on a different random stream, cut the arm at negative $x_0$ as well and came out
31\% too narrow in $x_0$.

\begin{cautionbox}[Reading hybrid widths]
Treat widths and tail quantities from a hybrid run as unconfirmed unless an independent expensive
chain checks them. Means are far more robust than tails.
\end{cautionbox}

\subsection{What the run writes}

For convenience, the run writes \code{work_dir/hybrid_diagnostics.txt} automatically. It is
rewritten in full every round, so a killed run still leaves its complete history and a resumed run
continues the same table. It is reproduced here in part only to explain its content. It opens with
the run's parameters and its evaluation budget:

\begin{lstlisting}[style=plain, basicstyle=\ttfamily\footnotesize]
expensive evaluations
  round 0 MCMC: (num_expensive_iters + 1) * num_walkers = (10 + 1) * 30 = 330
  regularization points: 50
  => raw training points at the first surrogate fit = 330 + 50 = 380
  each refinement round then adds num_validation_points = 100
  (exact unless a proposed point repeats one already evaluated, or an evaluation
   fails, in which case these are upper bounds)
\end{lstlisting}

Each verification spends one more validation batch, which is why this run's total is
$380 + 7 \times 100 + 3 \times 100 = 1\,380$.

A legend describing every column follows, and then the per-round table itself. Here is this run's,
with four columns dropped to fit the page: \code{surrogate}, which names the class fitted that round
and reads \code{GaussianProcess} throughout here, and the timings \code{t_valid_s},
\code{t_refresh_s} and \code{t_round_s} --- the validation batch, the expensive refresh, and the
round total.

\begin{lstlisting}[style=plain, basicstyle=\ttfamily\scriptsize]
round n_expensive n_train w_log_err post_shift surr_iters  Lc/tau  ess/tau  t_fit_s t_mcmc_s
--------------------------------------------------------------------------------------------
    1         380     123    0.7636        inf       5000    20.3      130      2.9     14.3
    2         480     223    0.0179      0.867       5000    18.9      121      2.4     16.4
    3         580     323    0.0078      0.043       5000    19.0      122      7.6     18.5
   3v                        0.0239      0.153      20000    51.6      330              81.0
    4         780     523    0.0098      0.224       5000    18.4      118     20.1     19.4
   4v                        0.0264      0.149      40000    59.4      380             190.6
    5         980     723    0.0280      0.101       5000    17.1      109     42.0     26.1
    6        1080     823    0.0102      0.243       5000    18.1      116     56.9     28.0
    7        1180     923    0.0035      0.117       5000    14.8       95     65.8     30.8
   7v                        0.0036      0.112      30000    50.6      324             226.4
\end{lstlisting}

Reading across a row:

\begin{itemize}
  \item \code{n_expensive} and \code{n_train} are the trim: what has been evaluated, and what
        entered the fit.
  \item \code{w_log_err} and \code{post_shift} are the two criteria, $\varepsilon$ in nats and
        $\delta$ in posterior standard deviations. They are what the run stops on --- here on
        \code{w_log_err} alone, since \code{posterior_shift_tolerance} is \code{None}, but
        \code{post_shift} is computed and printed regardless, so it can be read after the fact.
        The \code{inf} of round 1 marks that there is no previous posterior to compare against yet.
  \item \code{surr_iters}, \code{Lc/tau} and \code{ess/tau} are how deeply the cheap chain was
        sampled: the length actually run, and what it bought in autocorrelation times and effective
        samples. The refinement rows sit at the 5\,000-iteration minimum; the verification rows are
        where \code{surr_iters} jumps to 20\,000--40\,000 until both floors are cleared. A run that
        ends because its chains keep growing without clearing them shows it in these columns.
  \item \code{t_fit_s} against \code{t_mcmc_s} is the pair worth watching. The fit grows with the
        training set, 2.4\,s on 223 points to 65.8\,s on 923, while the cheap MCMC grows with the
        chain length asked of it: the 226.4\,s of round 7's verification is a 30\,000-iteration
        chain on the largest surrogate. On a long run it is the fits, not the sampling, that come to
        dominate --- which is when \code{max_train_points} or a \code{DistributedGPSurrogate}
        becomes worthwhile.
  \item A verified round prints two rows, $n$ and $n$v, and the columns they share ---
        \code{n_expensive}, \code{n_train}, \code{surrogate} and \code{t_fit_s} --- are left empty
        on the second: one fit, one training set, scored twice.
\end{itemize}

\section{Restarting, extending, and changing surrogate}

With \code{save_restart=True}, the restart file is written atomically after round 0 and after every
round, so a crash costs at most the round in progress. It carries the whole training set, which also
makes it the cleanest way to \emph{extend} a finished run: raise \code{num_rounds_max} and resume
with \code{load_restart=True}.

The restart file stores only surrogate-agnostic state --- the training points and their values, the
evaluation count, the per-round diagnostics and the walker positions --- so the surrogate is rebuilt
from whatever is passed on resuming. A run whose fits have become too slow can be stopped and resumed
as an ensemble on the same training data:
\begin{lstlisting}
# start with the dense Gaussian process, for its accuracy
slurm_mcmc_hybrid(..., surrogate='gp', save_restart=True, work_dir='run17')

# fits too slow? stop it, and resume as an ensemble on the same data
from slurmcmc.hybrid import DistributedGPSurrogate
slurm_mcmc_hybrid(..., load_restart=True, save_restart=True, work_dir='run17',
                  num_rounds_max=40,
                  surrogate=DistributedGPSurrogate(max_points_per_expert=500,
                                                   parallel='joblib'))
\end{lstlisting}
Every combination works, in either direction, and the switch is recorded per round. With
\code{save_surrogate}, fitted surrogates are written to \code{work_dir/surrogates/round{k}.pkl}:
\code{'latest'} keeps only the most recent, enough for a resumed round to skip its refit;
\code{'all'} keeps every round, at roughly $8n^2$ bytes each. A saved fit is reused only when the
round, the training-set size and the surrogate class all match.

Either mode also appends one line per fit to \code{work_dir/surrogates/surrogate_log.txt}: the
round, the surrogate class, how many points were accumulated, trimmed, subsampled and fitted, how
long the fit took, and the settings it was built with together with its fitted kernel. Being
appended rather than rewritten, it outlives \code{'latest'} pruning the pickles and continues
across a restart, so it is where a switch of surrogate --- or a fit that is becoming the bottleneck
--- shows up as a sequence.

\section{Results and warnings}

The returned dictionary contains the final posterior \code{samples}; \code{converged},
\code{num_rounds} and \code{num_expensive_evals}; per-round lists of every quantity that decided the
run, including \code{weighted_log_error_per_round}, \code{posterior_displacement_per_round},
\code{surrogate_iters_per_round}, \code{surrogate_tau_per_round}, \code{surrogate_ess_per_round},
\code{verified_per_round} and \code{blocking_criteria_per_round}, which names the criterion that held
each round back; \code{round_records}, one record per measurement including both rows of a verified
round; the fitted \code{surrogate}; the \code{validation_points} and \code{log_importance_weights}
of the final validation; and timings. The emulator's own uncertainty on the validation points,
\code{surrogate_std_per_round}, lets the common assumption that it is negligible be checked rather
than assumed.

Warnings are logged when refinement is not improving fast enough to be worth continuing; when
posterior mass piles up against \code{param_bounds} or the confinement envelope, which can signal a
model error, an unidentifiable parameter or bounds set too tightly; and when a surrogate chain hits its
iteration cap short of its criteria. The stall warning is tuned on \code{HybridMCMCConfig} by
\code{stall_relative_improvement}, \code{stall_window} and \code{stall_significance}.

\section{Argument reference}
\label{sec:hybrid-args}

Arguments shared with the other functions are in \cref{app:common}.

\subsubsection{Starting the run and its expensive budget}
\begin{argtable}
  \code{log_prob_fun} & required & the expensive log-probability; a callable or a deferred-import dictionary \\
  \code{init_points} & \code{None} & starting walkers of the round-0 MCMC, shape $(\Nc, d)$; required unless \code{initial_train_points} is given, and whenever \code{num_expensive_iters_per_round} $> 0$ \\
  \code{param_bounds} & \code{None} & a hard box that constrains the posterior; \code{None} runs unbounded (\cref{sec:unbounded}) \\
  \code{constraint_fun} & \code{None} & feasible when $\le 0$ \\
  \code{num_expensive_iters} & \code{20} & length of the round-0 MCMC; with the walkers, it sets the initial training set, $(\text{iters}+1)\times\Nc$ points \\
  \code{initial_train_points}, \code{initial_train_values} & \code{None} & existing evaluations, shapes $(n, d)$ and $(n)$; round 0 is skipped \\
  \code{num_regularization_points} & \code{0} & points placed without reference to a posterior: uniform in the box, or around the training data when unbounded \\
  \code{num_validation_points} & \code{100} & expensive evaluations per round; sets both the precision of the test and the cost of a round \\
  \code{num_rounds_max} & \code{5} & safety cap on the number of rounds \\
  \code{num_expensive_iters_per_round} & \code{0} & optional expensive MCMC each round, started from the surrogate posterior, costing $(\text{iters}+1)\times\Nc$ evaluations \\
\end{argtable}

\subsubsection{Surrogate and its sampling}
\begin{argtable}
  \code{surrogate} & \code{'gp'} & \code{'gp'}, \code{'polynomial'}, \code{DistributedGPSurrogate(...)}, or any object with \code{fit} and \code{predict} \\
  \code{polynomial_degree} & \code{3} & for \code{surrogate='polynomial'} \\
  \code{surrogate_uncertainty_penalty} & \code{1.0} & $\kappa$ in $\mu - \kappa\sigma$; \code{0} disables \\
  \code{num_surrogate_iters} & \code{2000} & block size of the surrogate chain \\
  \code{max_surrogate_iters_multiplier} & \code{8} & cap on the chain's extension, in blocks \\
  \code{num_surrogate_walkers} & \code{None} & walkers of the surrogate chain; \code{None} is \texttt{emcee}'s minimum, $2d+2$ \\
  \code{surrogate_burnin_fraction} & \code{0.2} & fraction of each chain discarded \\
  \code{refine_surrogate_ess}, \code{final_surrogate_ess} & \code{None}, \code{None} & ESS floors of the two depths; $10\,d$ and $100\,d$ \\
  \code{refine_iters_per_tau}, \code{final_iters_per_tau} & \code{None}, \code{50} & $\Lc/\tau$ floors of the two depths \\
  \code{max_verification_attempts} & \code{None} & rejections allowed before the run gives up \\
\end{argtable}

\subsubsection{Convergence and training data}
\begin{argtable}
  \code{log_error_threshold} & \code{0.1} & $\varepsilon^\ast$ in nats; tighten to 0.01--0.02 for tail quantities \\
  \code{min_ess_weights} & \code{10.0} & the degeneracy guard \\
  \code{posterior_shift_tolerance} & \code{0.1} & $\delta^\ast$ in posterior standard deviations; \code{None} disables \\
  \code{train_log_prob_trim_range} & \code{100.0} & trim points this many nats below the best; \code{None} disables \\
  \code{min_train_points_after_trim} & \code{50} & floor that stops the trim on small training sets \\
  \code{max_train_points} & \code{None} & cap on the fit, by uniform random subsample \\
  \code{train_log_prob_floor} & \code{None} & absolute floor on training values \\
\end{argtable}

\subsubsection{Saving}
\begin{argtable}
  \code{save_posterior_samples} & \code{None} & draws of each round's posterior kept in \code{work_dir/posteriors} \\
  \code{save_surrogate} & \code{None} & \code{'latest'} or \code{'all'} fitted surrogates in \code{work_dir/surrogates} \\
  \code{restart_file} & \code{'hybrid_restart.pkl'} & name of the restart file in \code{work_dir} \\
  \code{expensive_mcmc_kwargs} & \code{None} & forwarded to the expensive \code{slurm_mcmc} runs \\
\end{argtable}
